% !PRINTED 100
\weberpage{100}
%<<CONT
\[ \begin{aligned}
m = 819:\quad &\eta^3 - 273\eta + 91 = 0, & (\eta = \xi) \\
&\eta^3 - 273\eta + 91 . 19 = 0, \\
&\eta^3 - 273\eta + 91 . 8 = 0, \\
&\eta^3 - 273\eta + 91 . 17 = 0.
\end{aligned} \]

Die Wurzeln dieser einzelnen Gleichungen findet man am besten aus der
Gleichung (9) oder (11), indem man darin nach Potenzen von $\varrho$
ordnet, dabei aber nur $\varrho^3 = 1$, nicht $\varrho^2 + \varrho + 1 = 0$
benutzt.

So bekommt man z.\,B. für den Fall $m = 63$:

\[ \begin{aligned}
(\varrho, \eta) &= (\eta_1 + \varrho\eta_1' + \varrho^2\eta_1'')(\eta_2 + \varrho\eta_2' + \varrho^2\eta_2'') \\
\text{oder} \\
&= (\eta_1 + \varrho\eta_1' + \varrho^2\eta_1'')(\eta_2 + \varrho\eta_2'' + \varrho^2\eta_2'),
\end{aligned} \]

worin

\[ \begin{aligned}
\eta_1 &= r_1 + r_1^{-1}, & \eta_1' &= r_1^2 + r_1^{-2}, & \eta_1'' &= r_1^4 + r_1^{-4} \\
\eta_2 &= r_2 + r_2^{-1}, & \eta_2' &= r_2^3 + r_2^{-3}, & \eta_2'' &= r_2^2 + r_2^{-2}
\end{aligned} \]

zu setzen ist (Bd.\,I, \S.\,172), und man findet so:

\[ \begin{aligned}
\eta &= \eta_1\eta_2 + \eta_1'\eta_2'' + \eta_1''\eta_2' \\
\text{oder} \\
\eta &= \eta_1\eta_2 + \eta_1'\eta_2' + \eta_1''\eta_2''.
\end{aligned} \]

Um die Gruppen $\mathfrak{A}$ daraus zu finden, hat man zwei Zahlen $x, y$
so zu bestimmen, dass $7x + 9y = 1$ wird, und dann

\[ r = r_1 r_2, \quad \text{also} \quad r_1 = r^{7x}, \ r_2 = r^{9y} \]

zu setzen. Nimmt man $x = 4, y = -3$ an, so erhält man für $x$ die beiden % [sic: Weber prints ``für x'', although x has just been fixed and values of \eta follow.]
Werthe von $\eta$:

\[ \begin{aligned}
\eta &= (r^{28} + r^{-28})(r^{27} + r^{-27}) + (r^7 + r^{-7})(r^9 + r^{-9}) \\
&\qquad + (r^{14} + r^{-14})(r^{18} + r^{-18}) \\
\eta &= (r^{28} + r^{-28})(r^{27} + r^{-27}) + (r^7 + r^{-7})(r^{18} + r^{-18}) \\
&\qquad + (r^{14} + r^{-14})(r^9 + r^{-9}),
\end{aligned} \]

also

\[ \mathfrak{A} = \pm 1, \ \pm 8, \ \pm 2, \ \pm 16, \ \pm 4, \ \pm 32, \]

oder

\[ \mathfrak{A} = \pm 1, \ \pm 8, \ \pm 11, \ \pm 25, \ \pm 5, \ \pm 23. \]

Die beiden Gruppen können durch die Potenzen von 2 und von 11 dargestellt
werden.

Die cubischen Gleichungen, deren Bildungsgesetze wir jetzt kennen gelernt
haben, enthalten mit ihren \gesperrt{Tschirnhausen-Transformationen alle
cubischen Kreistheilungsgleichungen.} Aber diese Gleichungen sind auch alle
von einander
%CONT>>
